\documentclass[11pt,a4paper]{article}
\usepackage[T1]{fontenc}
\usepackage[utf8]{inputenc}
\usepackage{lmodern}
\usepackage[a4paper,margin=1in]{geometry}
\usepackage{microtype}
\usepackage{enumitem}
\usepackage{amsmath,amssymb}
\usepackage{mathtools}
\usepackage{hyperref}
\hypersetup{
	colorlinks=true,
	linkcolor=black,
	urlcolor=blue,
	citecolor=black
}

\setlist[itemize]{leftmargin=1.2em}
\setlist[enumerate]{leftmargin=1.2em}

\newcommand{\Title}[1]{\section*{#1}\vspace{-0.25em}\hrule\vspace{0.8em}}
\newcommand{\Field}[2]{\textbf{#1:}~#2\par}
\newcommand{\Affil}[1]{\textit{#1}\par}
\newcommand{\Abstract}{\vspace{0.6em}\textbf{Abstract.}~}

\begin{document}
	
	\begin{center}
		{\LARGE Online Algebra, Logic, and Topology Seminar}\\[0.25em]
		{\large Category Theory Seminar — Abstract}\\[0.5em]
		December 02nd, 2025
	\end{center}
	
	\vspace{1.25em}\hrule\vspace{1.25em}
	
	\Title{Categorical and metric density}
	
	\Field{Speaker}{Tome Leinster}
	\Affil{University of Edinburgh, UK}
	
	\Field{Date \& Time}{December 02nd, 2025 (Tuesday) — 3{:}00\,PM (Coimbra time)}
	\Field{Place}{Google Meet}
	
	\Abstract
When metric spaces are viewed as enriched categories, the categorical
notion of density does not coincide with the ordinary topological notion.
This prompts two questions. First, what does categorical density mean for
metric spaces? And second, does the condition of topological density
generalize from metric spaces to arbitrary enriched categories? Both
questions were considered in Lawvere's 1973 metric spaces paper, but only
very briefly: there is a lot more to say. Some satisfying new answers have
recently been uncovered by my student Adrián Doña Mateo, which I will
explain.	
\end{document}
